\section{问题分析}

四个问题的物理内核是同一个：热风与药材表面之间的对流换热、药材内部的导热，以及水分在浓度梯度驱动下的内部扩散与表面传质。四问的差别只在于时段长度、物性是否随状态变化、是否需要判定终点、以及边界是否会动。因此本节按``主导机制、难点、模型应做到什么程度三条线索分别分析，为第~\ref{sec:assumptions} 节的假设取舍与第~\ref{sec:modelbase} 节的统一建模提供依据。

\subsection{问题一：预热阶段的时间尺度与两个过程的解耦}

问题一的关键判断有两个。其一是\textbf{温度场在本问时段内远未平衡}。由附录 2 的常数物性可算出热扩散率与导热特征时间
\begin{equation}
\alpha=\frac{\kcond}{\rho\cp}=\frac{0.36}{820\times2600}=1.689\times10^{-7}~\mathrm{m^2/s},
\qquad
\tau_T=\frac{{\Rzero}^2}{\alpha}=2369~\mathrm{s}=39.5~\mathrm{min},
\label{eq:tauT}
\end{equation}
即题面给出的 $1800$~s 只相当于 $0.76\,\tau_T$。这说明预热阶段药材内部存在明显的径向温度梯度，不能采用``药材整体同温''的集中参数假设。其二是\textbf{水分过程的特征时间比传热高两个数量级}：附录 2 的扩散系数为 $\Ddiff=7\times10^{-9}e^{-0.89/\C}$，在初值处
\begin{equation}
\Ddiff({\C}_0)=4.938\times10^{-9}~\mathrm{m^2/s},
\qquad
\tau_C=\frac{{\Rzero}^2}{\Ddiff({\C}_0)}=8.10\times10^{4}~\mathrm{s}=22.5~\mathrm{h},
\label{eq:tauC}
\end{equation}
$t_{\mathrm{end}}/\tau_C$ 仅 $2.2\%$。由此可以预期：全场平均含水率几乎不变，但\textbf{近表面薄层}的水分变化不可忽略——扩散穿透深度为 $\sqrt{\Ddiff t_{\mathrm{end}}}=2.98\times10^{-3}$~m，即 $0.15\Rzero$，恰好落在题目要求输出的 $0\sim2$~cm 网格的表层几列上。这两个判断共同决定了模型既不能退化为``只解热传导、含水率视为常数''，也不能对水分场做准稳态近似。

此外，附录 2 中 $\rho$、$\cp$、$\kcond$ 均为常数、$\Ddiff$ 只依赖含水率而不依赖温度，因此\textbf{温度方程与水分方程在本问完全解耦}，可以分别求解，这既降低了实现复杂度，也使误差来源可分离（第~\ref{sec:validation} 节据此把两者分别与解析解对照）。两个无量纲阻力比也说明两侧边界都不可忽略：
\begin{equation}
\frac{\hconv\Rzero}{\kcond}=\frac{25\times0.02}{0.36}=1.39,
\qquad
\frac{\hmass\Rzero}{\Ddiff({\C}_0)}=\frac{8\times10^{-7}\times0.02}{4.938\times10^{-9}}=3.24,
\end{equation}
二者均为 $O(1)$，即表面换热阻力与内部导热阻力同量级、表面传质阻力与内部扩散阻力同量级，故表面的 Robin 型边界必须显式处理，不能简化为给定表面温度或给定表面含水率。

\subsection{问题二：变物性带来的双向耦合与时间尺度分离}

问题二在物理机制上与问题一的本质区别是\textbf{物性不再是常数}。按题意统一采用附录 3 后：$\rho$ 依赖含水率，$\cp$ 与 $\kcond$ 依赖湿基质量分数 $\wfrac=\C/(1+\C)$，而 $\Ddiff$ 同时依赖含水率与\textbf{温度}（$\Ddiff=2.4\times10^{-3}e^{-0.45/\C}e^{-3850/\T}$，其中 $\T$ 必须用开尔文）。这带来三重后果。

第一，\textbf{耦合方向是不对称的}。温度通过指数因子 $e^{-3850/\T}$ 影响扩散系数：温度由 $301.15$~K 升到 $323.315$~K 时该因子为
\begin{equation}
\exp\!\left[3850\left(\frac{1}{301.15}-\frac{1}{323.315}\right)\right]=2.39,
\end{equation}
即升温 $22$~K 使 $\Ddiff$ 增大近 $2.4$ 倍；而含水率的下降对 $\Ddiff$ 的影响要大得多：由初值 $2.55$ 降到判据值 $0.15$ 时，$e^{-0.45/\C}$ 由 $0.705$ 降到 $0.0498$，$\Ddiff$ 由 $5.64\times10^{-9}$ 降到 $8.05\times10^{-10}~\mathrm{m^2/s}$，\textbf{下降约 7 倍}；若再降到表面实际的 $0.05$ 量级，$\Ddiff$ 将降至 $10^{-12}$ 量级。也就是说，干燥后期内部水分向外补充的能力会急剧恶化。

第二，\textbf{时间的快慢由两个尺度分别决定}。初值处导热特征时间 ${\Rzero}^2\rho\cp/\kcond=0.77$~h，远小于 $2\sim3$ 天的干燥时长，故温度场会很快趋于均匀并逼近环境温度；而扩散特征时间由 $19.7$~h（初值）增大到 $138.1$~h（判据值处），与干燥总时长同量级甚至更大。因此\textbf{整个过程的时间尺度由传质决定}。这一点还可以用表面传质阻力与内部扩散阻力之比确认：该比值由初值处的 $2.84$ 升到后期的 $20.5$，说明后期表面水分被气流迅速带走、而内部供给跟不上，属于典型的\textbf{内部扩散控制}；其直接后果是干燥后期表面附近会形成很陡的浓度边界层，这正是本文采用表面加密网格的动机。

第三，\textbf{两阶段在数学上不需要分开处理}。题目问题 2 明确``为简化问题，相关经验公式统一采用附录 3''，即两阶段的参数差异被题面显式抹平，所剩下的差别只有环境条件 $T_\infty(t)$、$\Cinf(t)$ 的时程。而附件 1 只覆盖 $0\sim4$~h，其后必须做外推假设——这是本题最大的不确定性来源（第~\ref{sec:sensitivity} 节量化）。因此在建模上，本文把``两阶段''理解为环境驱动的两个不同时段，而不引入任何人为的阶段切换判据，避免凭空增加一个题目未给出的参数。

\subsection{问题三：判据的数学表述与求根口径}

问题三在方程层面与问题二完全相同（同一物性集、同一几何、同一初值），新增的只有终点判据。因此本问的难点全在``如何把题目语言翻译成一个可计算的数学问题''。

\textbf{其一是``各处''的解释。} 题目要求药材各处的水分浓度低于 $0.15$~kg/kg。由于水分由内向外扩散，含水率沿半径单调递减，最慢达标的位置是\textbf{轴心}，故``各处达标''等价于全场最大值达标：
\begin{equation}
M(\t):=\max_{0\le\rpos\le\Rzero}\C(\rpos,\t)\le 0.15~\mathrm{kg/kg}.
\label{eq:maxcrit}
\end{equation}
这是所有解释中最保守的一种（若按断面平均值或表面值理解，烘干时间会更短）。算例实测也印证了这一点：交付结果文件中每一行的最大值恒出现在第 0 列（即 $\rpos=0$）。

\textbf{其二是严格不等号与舍入的关系。} 判据是严格不等式，而交付文件按题目要求只保留四位小数，于是存在两个相邻的采样行都显示 $0.1500$ 的情形：仅凭交付表无法判断哪一行才真正达标。因此本文把判据定义在\textbf{内部求解网格的未舍入值}上，并把``达标时刻''作为数学问题的一个部分来求解，而不是从交付表反推。同时明确区分两个口径：内部连续过零时刻 $\tend$（判据严格成立的时刻）与交付网格上首个严格合格的时刻（$(\lfloor\tend/\dt_{\mathrm{out}}\rfloor+1)\dt_{\mathrm{out}}$，保证交付文件只含等间隔行）。两者都给，但不混用。

\textbf{其三是计算量与算法选择。} 本问要在 $2\sim3$ 天的时间尺度上确定一个阈值时刻。若用二分法，每次迭代都要重跑全程，代价随精度要求线性增长；而 $M(\t)$ 是有限个连续函数取最大值，因而连续但只在 $\arg\max$ 切换处不可微。这一性质正好适合\textbf{过零事件检测}：只需一次积分，让积分器在检测到 $M(\t)-0.15$ 变号时自动定位并终止，即可在同一精度下得到 $\tend$。因此本问的求解策略是``单次积分 $+$ 事件检测''，而不是二分或粗网格扫描。

\subsection{问题四：移动边界、两个相反效应与参照算例的必要性}

问题四同时改变了两件事：物性公式由附录 3 换成附录 4，几何由固定半径换成附件 2 给出的收缩半径。此外，附件 2 只给出半径，未给长度与体积，因此``失水如何改变形状''本身就需要额外的建模假设。这些使得问题四成为四问中建模工作量最大、也最容易得出错误定性结论的一问。

\textbf{其一是移动边界。} 求解域本身随时间缩小，直接后果是空间离散的网格要么随之运动（ALE/动网格），要么把方程变换到固定域上求解。后者只需引入材料坐标 $\xi=\rpos/\R(\t)$，实现代价低一个数量级，且在本问的``径向均匀收缩''假设下，固定 $\xi$ 的点恰好始终跟踪同一团材料，$\xi$ 就是材料标签。因此本文选择坐标变换路线。

\textbf{其二是两个方向相反的效应。} 直觉上，半径减小一方面使水分的逸出面积变小（不利于干燥），另一方面使内部水分的扩散路径变短（有利于干燥），两个效应方向相反。在本题采用的边界闭合下哪一个占优，必须靠推导确定而不能凭直觉。本文在材料坐标下对守恒方程沿断面积分，得到单位材料测度的失水速率与半径成反比、单位长度药材的失水速率正比于 $1/\R$，而内部再分布的时间尺度正比于 ${\R}^2/\Ddiff$：两条通道都指向``变快''，净效应的量级应与 $({\R}_{\mathrm{end}}/\Rzero)^2$ 相当。这一推论随后被对照算例定量验证（第~\ref{sec:q4result} 节）。

\textbf{其三是参照算例的必要性。} 问题三与问题四的答案分别为 $57.17$~h 与 $50.82$~h，二者相差 $11\%$，但这两问\textbf{同时}改变了物性集与几何。因此这一差值\textbf{不能}被解释为``收缩使干燥变快''。要分离两个效应，必须构造一个\textbf{同物性集、只变几何}的对照算例（附录 4 物性 $+$ 固定半径），用``对照算例 $\to$ 问题四''这一对差值度量收缩效应、用``问题三 $\to$ 对照算例''这一对差值度量参数集效应。本文按此口径给出三算例对照，并在第~\ref{sec:q4result} 节说明其局限。这一处理是问题四结论可靠性的关键，也是本文与``直接比较 $57.17$ 与 $50.82$''的做法之间最重要的区别。

\subsection{四问的难点与对策汇总}

